\begin{claim}{C.67}{Proved}{The exact relative disk group for a locally connected frontier}
Let $O\subset\R^2$ be a nonempty bounded simply
connected open set with locally connected frontier.
Choose a Riemann map $\phi:\mathbb D\to O$ and
retain its continuous Carath\'eodory extension
\[
                  q:\overline{\mathbb D}\longrightarrow\overline O.
                                                               \tag{C.67.1}
\]
Put
\[
 \begin{split}
 \mathcal D_\partial
     &=\{a\in\operatorname{Homeo}(\overline{\mathbb D}):
                                      a|_{S^1}=\id\},\\
 \mathcal K_O
     &=\{b\in\operatorname{Homeo}(\R^2):
                                b|_{\R^2\setminus O}=\id\}.
 \end{split}                                             \tag{C.67.2}
\]
The actual map
\[
 T_q:\mathcal D_\partial\longrightarrow\mathcal K_O,
 \qquad
 T_q(a)|_O=\phi a\phi^{-1},\qquad
 T_q(a)|_{\R^2\setminus O}=\id
                                                               \tag{C.67.3}
\]
is an isomorphism of topological groups with their
compact-open and inverse topologies.
In particular $\mathcal K_O$ is contractible.
The frontier need not be a Jordan curve.

For the actual stabilizer $H_O$ of $O$ in any
specified group of ambient operators, let
\[
 \eta_q:H_O\longrightarrow
                   \operatorname{Homeo}(\overline{\mathbb D})
                                                               \tag{C.67.4}
\]
be its original prime-end action in the chosen
Riemann coordinate. Its kernel is exactly
$\{h\in H_O:h|_O=\id\}$. The isomorphism (C.67.3)
restricts to the exact comparison
\[
 C_{\mathcal D_\partial}(\eta_q(H_O))
   \xrightarrow{\ \cong\ }
       \mathcal K_O\cap
               C_{\operatorname{Homeo}(\R^2)}(H_O).
                                                               \tag{C.67.5}
\]
Both directions have the actual induced group
topologies. The contraction asserted above is
for the full relative group $\mathcal K_O$;
no equivariance of the Alexander contraction
under a nontrivial operator is asserted.
\end{claim}

\begin{proof}
\emph{The retained quotient and prime-end maps.}
The Carath\'eodory theorem for locally connected
boundaries gives (C.67.1); it is a continuous
surjection, equals $\phi$ on the open disk,
and maps the boundary circle onto $\partial O$.
The same conformal coordinate identifies the
prime-end compactification of $O$ with the
closed disk. We use precisely the crosscut-chain
compactification and the extension theorem
stated in \cite[Sections~1 and~3.1]
{C_BinderRojasYampolsky2014}.
The map $q$ may have several preimages for a
frontier point. No inverse of $q$ on the
frontier is used.

An ambient homeomorphism $h$ preserving $O$
induces a homeomorphism of this prime-end
compactification. Here is the relevant
functorial construction. It sends a crosscut,
with its two distinct frontier endpoints,
to the crosscut with their actual images.
It sends the corresponding complementary
crosscut regions to their actual images.
On the compact set $\overline O$, $h$ is
uniformly continuous. Thus disjoint crosscut
closures remain disjoint and diameters
tending to zero still tend to zero.
Nested fundamental chains are sent to nested
fundamental chains. The inclusion relations
defining equivalence of chains are preserved,
as are those relations under $h^{-1}$.
The image-chain rule is therefore a bijection
on prime ends. Basic crosscut-region
neighborhoods and the inverse rule show that
its extension is a homeomorphism.
This construction is independent of the
auxiliary interior basepoint used to select
crosscut sides: a chain tends to the frontier
and eventually excludes either chosen
interior point.

Conjugating this extension by the fixed
prime-end Riemann coordinate gives
$\widehat h=\eta_q(h)$, with the exact equations
\[
 \widehat h|_{\mathbb D}=\phi^{-1}h\phi,\qquad
                    q\widehat h=hq .
                                                               \tag{C.67.6}
\]
The second equation holds first on the
dense open disk and then everywhere by
continuity. Composition of image chains,
or uniqueness on the dense disk, proves
that $\eta_q$ is a homomorphism.
It is identity exactly when $h|_O$ is
identity. These are the asserted original
restriction labels and kernel.

\emph{A frontier-fixing map fixes the entire prime-end circle.}
Let $b\in\mathcal K_O$ and let $\widehat b$
be its just constructed prime-end extension.
It preserves orientation: $b$ is identity
on the nonempty open exterior of the
bounded set $\overline O$, so its global
orientation is positive.
Write $r=\widehat b|_{S^1}$ and
$q_\partial=q|_{S^1}$. Since $b$ fixes
$\partial O$ pointwise, (C.67.6) gives
\[
                         q_\partial r=q_\partial.
                                                               \tag{C.67.7}
\]
We prove that this exact equation forces
$r=\id$.

First, $q_\partial$ is constant on no
nonempty open arc. Otherwise, for a point
$p$ equal to that constant value, the
holomorphic map $\phi-p$ is continuous
up to an open circle arc and is zero
there. In a local conformal coordinate
straightening that arc to a real segment,
Schwarz reflection extends it holomorphically
across the segment. The identity theorem
then makes $\phi-p$ identically zero,
contrary to injectivity of $\phi$.

Fix $x_*\in O$. The circle-valued map
\[
 v(\zeta)=\frac{q_\partial(\zeta)-x_*}
                    {|q_\partial(\zeta)-x_*|}
                  \quad(\zeta\in S^1)                 \tag{C.67.8}
\]
has degree one. Indeed $\phi(sS^1)$,
for $s<1$ sufficiently close to one,
is a positively oriented Jordan curve
enclosing $x_*$. Uniform convergence of
$\phi(s\zeta)$ to $q_\partial(\zeta)$
and the positive distance from $x_*$
to $\partial O$ preserve that winding
number in the limit.

Use $\zeta=e^{2\pi i t}$ and choose
a real lift $R$ of $v$, so
$R(t+1)=R(t)+1$. For an increasing
lift $P$ of $r$, equation (C.67.7)
implies that $R(P(t))-R(t)$ is one
constant integer. Replacing $P$ by
the unique integer translate which
subtracts that integer gives
\[
                    R(P(t))=R(t)\quad(t\in\R).
                                                               \tag{C.67.9}
\]
The function $R(t)-t$ is periodic and
bounded. Thus the full orbit
$\{P^n(t):n\in\Z\}$ is bounded for
each $t$, since its $R$-value is
constant. An increasing real
homeomorphism with such an orbit
has a fixed point: a nonconstant
forward monotone orbit has a finite
limit, and continuity fixes that
limit. Hence $P$, and therefore
$r$, has a fixed point.
Moreover, if $r(e^{2\pi i t})=e^{2\pi i t}$,
then $P(t)=t+k$ for an integer $k$.
Periodicity would give $P^n(t)=t+nk$;
the established boundedness forces $k=0$.
Thus the fixed set of $r$ is exactly
the projection of $\operatorname{Fix}(P)$.

If $r$ were nonidentity, choose
a component of the complement
of its fixed set. Using the
periodic fixed set of $P$, lift
that component to an interval
between consecutive boundary
points of $\operatorname{Fix}(P)$.
The interval is bounded, its
endpoints are fixed, and $P$
has no fixed point inside.
Iterations of every interior
point therefore tend to one
endpoint. Equation (C.67.7)
and continuity make $q_\partial$
constant on that open arc.
This contradicts the no-constant-arc
property. We conclude that
$r=\id$ as claimed.

\emph{The algebraic comparison and its topology.}
For $a\in\mathcal D_\partial$, the
continuous map $qa$ on the closed
disk is constant on every fiber
of $q$. Interior fibers are single
points; on a boundary fiber $a$
fixes every point. The compact
quotient $q$ therefore gives a
continuous map on $\overline O$,
equal to the identity on its
frontier. The identical construction
with $a^{-1}$ gives its inverse.
Pasting this map with identity on
the closed complement of $O$
defines the actual homeomorphism
(C.67.3). The formulas on $O$
give injectivity and preservation
of products.

Conversely, for $b\in\mathcal K_O$,
the prime-end extension $\widehat b$
fixes $S^1$ by the preceding argument.
Thus $\widehat b\in\mathcal D_\partial$,
and its interior formula gives
$T_q(\widehat b)=b$ everywhere.
This proves bijectivity.

For continuity let $\omega_q$ be
the modulus of uniform continuity
of $q$ on the compact closed disk.
The actual forward maps satisfy
\[
 \sup_{x\in\R^2}|T_q(a)(x)-T_q(a')(x)|
       \leq
       \omega_q\!\left(
        \sup_{\zeta\in\overline{\mathbb D}}
                                      |a(\zeta)-a'(\zeta)|
                    \right).                           \tag{C.67.10}
\]
Outside $O$ both sides of the
comparison of maps are identity.
The same estimate applies to
their inverses. Hence $T_q$
is continuous in the actual
group topologies.

The source $\mathcal D_\partial$
is a closed subgroup of the
Polish homeomorphism group of
the compact disk. The target
$\mathcal K_O$ is a closed
subgroup of the Polish plane
homeomorphism group: it is
defined by the pointwise
equations on $\R^2\setminus O$.
The source acts continuously
and transitively on this Polish
target by $(a,b)\mapsto T_q(a)b$.
The Polish open orbit-map
theorem verified in C.26 and
used in C.30 makes its orbit
map at identity open.
That orbit map is exactly
$T_q$. The continuous bijection
is therefore open and has a
continuous inverse. This proves
the full topological group
isomorphism, without assuming
injectivity of the boundary
quotient $q$.

\emph{Contraction of the full relative group.}
The Alexander family on
$\mathcal D_\partial$ is
\[
 A_s(a)(\zeta)=
 \begin{cases}
       s\,a(\zeta/s),&|\zeta|\leq s,\quad s>0,\\
       \zeta,&|\zeta|\geq s ,
 \end{cases}
 \qquad A_0(a)=\id.                                    \tag{C.67.11}
\]
The boundary-fixing condition
gives agreement on $|\zeta|=s$.
Every map has inverse
$A_s(a^{-1})$, and fixes $S^1$.
For $s>0$ joint uniform
continuity follows from uniform
continuity on the compact disk;
at zero both the map and its
inverse have displacement at
most $2s$. Also
$\|A_s(a)-A_s(a')\|_\infty
\leq s\|a-a'\|_\infty$.
These observations give joint
continuity in $s$ and $a$,
including inverses.
We have $A_1(a)=a$,
$A_0(a)=\id$, and
$A_s(\id)=\id$.
Transporting (C.67.11) by the
proved isomorphism contracts
the full group $\mathcal K_O$.

\emph{Every original operator relation.}
For $h\in H_O$, (C.67.6)
implies on the open disk,
and therefore by continuity
on the relative plane maps,
\[
 T_q(\widehat h\,a\,\widehat h^{-1})
       =h\,T_q(a)\,h^{-1}.
                                                               \tag{C.67.12}
\]
Outside $O$ both sides are
identity, since $h$ preserves
$O$ and its complement.
Conjugation by $\widehat h$
preserves $\mathcal D_\partial$
even when $h$ reverses orientation.
Equation (C.67.12) proves both
directions of the exact
centralizer comparison (C.67.5).
Restriction of a homeomorphism
of topological groups to these
subgroups retains their
induced topologies.
An original operator in the
kernel of $H_O\to\operatorname{Homeo}(O)$
has $\widehat h=\id$ and
conjugates every map of
$\mathcal K_O$ trivially;
that actual kernel has been
retained.

The same identity also keeps
the labels of operators which
transport $O$ to another such
domain $O'$. With a chosen
Riemann map $\phi'$ and extension
$q'$ there, the prime-end
transport is the unique disk
homeomorphism extending
$\phi'^{-1}h\phi$, and satisfies
$q'\widehat h=hq$.
Thus
\[
 T_{q'}(\widehat h\,a\,\widehat h^{-1})
        =h\,T_q(a)\,h^{-1}.                             \tag{C.67.13}
\]
Composing two original operators
composes these exact transports,
by their uniqueness on the
open disk.

Equations (C.67.5), (C.67.12),
and (C.67.13) provide the
operator-preserving disk models.
The Alexander family was proved
to contract the full relative
group. Its preservation of the
operator centralizers in
(C.67.5) has not been established.
Nor does this one-domain
topology comparison supply
uniform control for contractions
over infinitely many transported
domains. These are retained
extension requirements in C.12.
\end{proof}
